\scanpage{32}
% archive sheet number: 31

\deleted{$\nu$ \uline{premier à $p$} ($p$ fini) $\longleftrightarrow$}

% slanted
\margin{On suppose \uuline{$u$ \deleted{\unsure{$\ne$}}non central}}

% slanted
\margin{$\zeta$ une solution de $\zeta^2-\alpha\zeta+1=0$ dans $\bar k$}

$k$ \uline{car. 0} \uline{($p=+\infty$)}

\[
\begin{aligned}
\deleted{\unsure{A)|a)}}\quad \nu\ \text{fini}&\Longleftrightarrow\exists\,\deleted{\unsure{n}}\inserted{n\geq3}\ \text{tel que }\Psi_n(\alpha)={}\\ 
\deleted{\unsure{B)|b)}}\quad\text{''}&\Longleftrightarrow \deleted{\zeta\ \text{racine}}\ \zeta\ \text{d'ordre fini }\ne1,2
\end{aligned}
\]

$k$ car. $p>0$

\[
\begin{aligned}
\mathrm A)\quad\nu\ \text{fini}&\Longleftrightarrow\alpha\ \text{algébrique}\Longleftrightarrow n\geq1\ \text{tel}\\ 
&\hspace{4em}\text{que }\Psi_n(\alpha)=0\ (\\ 
&\Longleftrightarrow\exists\,n\geq1,\ \text{premier à }p\\ 
&\hspace{3em}\text{ou de la forme }\deleted{\unsure{2}}2p,\\ 
&\hspace{3em}\deleted{\unsure{\text{ou }2p\ (\text{avec}|2p\ (\text{avec}}}\ p\ \text{impair})\\ 
&\hspace{3em}\inserted{\text{ou }p,}\ \text{tel que }\Psi_n(\alpha)=0\\ 
(S\ne\varnothing)&\quad[\text{ledit }n\ \text{étant alors}\\ 
&\hspace{3em}\text{unique, et égal à }\nu(u)]\\ 
&\Longleftrightarrow\zeta\ \text{d'ordre fini. En fait,}\\ 
&\hspace{2em}\text{si}
\end{aligned}
\]

\uline{Théorème} $\widetilde G$ forme de $\mathrm{Sl}(2)_S$, $u$ section de $\widetilde G$, $\alpha=A(u)\in\Gamma(S,\unsure{\mathcal O_S|\frac{\mathcal O}{S}|\underline{O}_S})$,
$n\in\mathbb N$, $\boxed{n\geq3}$. On s'intéresse à la traduction, en termes de $\alpha$,
et de propriétés des car. résiduelles de $S$, de la propriété
suivante

% slanted
\margin{\uline{NB} On ne permet pas $n=2$, qui \unsure{compliquerait|cancellerait|se calculerait} complètement}

\[
(*)\quad u^n=\mathrm{id},\quad\text{et }\forall s\in S,\ \text{l'ordre de }u_s\ \text{est égal à }n.
\]

A) Supposons que $n$ ne soit \deleted{ni un nb premier, ni} \inserted{de la forme $p$ \inserted{ni $2p$} ($p$ premier}
\deleted{ni $2p$ ($p$ premier)} \deleted{impair} \uline{impair}). Alors $(*)$ équivaut à :

a) $\Psi_n(\alpha)=0$, et $n$ inv. sur $S$ (i.e. les car résiduelles
sont premières à $n$)

% slanted
\margin{\uline{NB} le cas $n=4$ \unsure{est pas|n'est pas} exclu -- A) doit marcher pareil ds le cas de $G=\widetilde G/\mu_2$}

B) $n=2p$ ($p$ premier impair). Alors $(*)$ équivaut à :

b) $\Psi_n(\alpha)=0$, et $2$ inv sur $S$ (i.e. car. rés. de $S$
sont $\ne2$)

% slanted
\margin{\uline{à vérifier} $\left.\vphantom{\Large A}\right\}$ Peut-être faux (\unsure{bon entraînement|bien certainement|bons contre-exemples} \unsure{faites sur papier|faux si non premier}…)}

C) $n=p$, $p$ est premier \inserted{impair}. Alors $(*)$ équivaut à :

c) $\Psi_n(\alpha)=0$.

\uline{NB} Ces conditions \uline{n'impliquent pas} $(\alpha+2)(\alpha-2)$ inv, sauf ds le cas A
Dans les cas B \deleted{resp. C} \inserted{et C}, les pts où $\alpha_s^2-4=0$ sont ceux de car $p$,
\deleted{(où $\alpha_s=-2\ne2$] (resp. ceux de car} où $\alpha_s=-2\ne+2$ (cas B) resp $\alpha_s=+2\ne-2$ (cas C)
